\SetAccent{E65B63}{FCEDEF}
\Hero{One-day applied learning program}{INNOVATION IN ACADEMIA}{A faculty transformation lab for stronger patents, funded research and industry relevance}{Research novelty \textbullet\ Patent potential \textbullet\ Fundable partnerships}{01}
\MetaBar{Teachers \& academic leaders}{1 day | 9:30-17:00}{Hands-on studio}{Discuss with us}
\vspace{3mm}
\WhyCard{Turn faculty expertise into a visible pipeline of novel research, defensible IP and fundable partnerships.}{Research scholars learn to identify consequential problems, build inventive confidence and see how research becomes IP and real-world impact.}{Faculty move beyond incremental publication, strengthen novelty and create credible paths to patents, funded research and industry partnerships.}
\CardGap
\TwinCard{What we will explore}{\BulletList{\item Build an inventive mindset for research, teaching and institutional challenges.\item Separate genuine novelty from routine extensions of existing work.\item Use structured frameworks to connect evidence, IP and stakeholder value.}}{What participants gain}{\BulletList{\item \textbf{Find novelty:} identify innovative, relevant problems worth solving.\item \textbf{Escape the publication trap:} move from incremental papers to defensible contribution.\item \textbf{Create IP:} shape novel, non-obvious and patent-ready concepts.\item \textbf{Attract support:} build credible industry and government-funding pathways.}}
\CardGap
\begin{DarkCard}{\color{BroGold}\bfseries\fontsize{8}{9.5}\selectfont INSTITUTIONAL RETURN\par}\vspace{1mm}{\color{white}\fontsize{10}{12}\selectfont\bfseries A stronger pipeline of patentable research, fundable proposals and industry-facing faculty initiatives.}\end{DarkCard}
\ProgramFooter{NOVEL PROBLEMS \textbullet\ DEFENSIBLE IDEAS \textbullet\ REAL-WORLD IMPACT}
\newpage
\CompactHero{Innovation in Academia}{From novelty to evidence, IP and partnership}{A clear one-day journey, followed by reusable tools and a practical continuation path.}
\JourneyTable{%
\JourneyTableRow{09:30-10:15}{Beyond the next paper}{Recognise the iterative publication trap and define meaningful impact.}
\JourneyTableRow{10:15-11:15}{Novel problem discovery}{Scan tensions, unmet needs, anomalies and emerging opportunity landscapes.}
\JourneyTableRow{11:30-12:45}{Storytelling for relevance}{Connect evidence, beneficiaries, urgency and a compelling research question.}
\JourneyTableRow{13:45-15:00}{Framework lab}{Use TRIZ contradictions, inventive principles and structured idea generation.}
\JourneyTableRow{15:15-16:15}{Reverse brainstorming and IP lens}{Expose failure modes and test novelty and non-obviousness.}
\JourneyTableRow{16:15-17:00}{Impact pathway}{Map experiments, patent steps, industry partners and suitable funding routes.}
}
\CardGap
\begin{minipage}[t]{.485\textwidth}
\SideCard{Who should attend}{\BulletList{\item University and college teachers\item Heads of department and deans\item Research supervisors and lab leaders\item Curriculum, accreditation and faculty-development teams}}
\CardGap
\SideCard{Takeaway toolkit}{\BulletList{\item A portfolio of consequential, underexplored research problems\item A novelty and non-obviousness canvas for a selected idea\item Concepts generated using TRIZ and reverse brainstorming\item A 30-day experiment, IP and stakeholder-engagement roadmap}}
\end{minipage}\hfill
\begin{minipage}[t]{.485\textwidth}
\SideCard{Methods \& tools}{\BulletList{\item Novel problem and opportunity landscapes\item Research storytelling for relevance\item TRIZ contradictions and principles\item Reverse brainstorming\item Publication-to-patent decision lens\item Industry collaboration and funding canvas}}
\CardGap
\CustomProgramCard{Custom programs can be discussed}{The scope can be shaped around the cohort, institutional priorities and desired outcomes. Workbook, innovation tool cards, IP lens, funding canvas and participation certificate can be included.}
\end{minipage}
\ProgramFooter{ONE COHORT \textbullet\ ONE CAPABILITY \textbullet\ REVIEW THE EVIDENCE}
\newpage
